\section{Methods and Benchmark Design}
\label{sec:methods}

\subsection{Experimental design and frozen model identity}

% Question: What is the experiment, and what exactly is the method under
% evaluation?
% Answer: A controlled recovery experiment on one method-neutral injected
% detector cube, with NOVA's frozen science release, NOVA-S Science R1,
% evaluated against extraction-based analyses of the same cube.
% Evidence: The NOVA-S handoff README and science contract.

The project uses a controlled recovery experiment. A detector-level injector
generates a synthetic SOSS time series with a known transmission spectrum. NOVA
and one or more extraction-based analyses receive the same detector cube, time
grid, uncertainty array, data-quality information, and declared calibration
products. Each method produces a fixed spectral prediction before the injected
spectrum is made available for scoring. The experiment therefore separates
model development from absolute evaluation.

NOVA is the general detector-level inference framework introduced in
Section~\ref{sec:structural-limitations-inference}. The specific
implementation evaluated in this assessment is its frozen science release,
NOVA-S Science R1, referred to below as NOVA-S. The release is immutable and
authenticated by content hashes; later response experiments, adaptive
candidates, and injector variants are diagnostics or prospective candidates
and are not part of NOVA-S. Throughout development, no injected truth or
recovered-depth residual may construct or select a response, mask, nuisance
model, threshold, or correction: every calibration uses only the observed
detector data, instrument reference files, and out-of-transit (OOT)
information, and the injected spectrum is read only after a detector
prediction has been closed. Model classes that failed truth-blind gates are
retained and reported rather than silently discarded
(Appendix~\ref{app:nova}).

\subsection{The V45 benchmark cube and retained detector data}

% Question: What data does the frozen benchmark contain?
% Answer: 229 Stage-2 rateints integrations built from real WASP-17
% out-of-transit data with a 127-integration synthetic event, consumed
% unchanged by every recovery method.
% Evidence: The exact NOVA-S method note and the injector evolution record.

The frozen benchmark is the V45 method-neutral cube: real out-of-transit
Stage-2 \texttt{rateints} integrations from the WASP-17 NIRISS/SOSS
observation provide the source-free temporal baseline, and a known
wavelength-dependent transit deficit, applied to a denoised per-order
expectation of the target source only, is added to it. Additive background
and unrelated contaminants therefore do not transit, and the realised noise
is not artificially modulated; the exact injector construction is given in
Appendix~\ref{app:benchmarking}. The delivered cube contains 229
integrations; the synthetic event occupies integrations 51--177, giving 127
event and 102 complement integrations. Stage-2 SCI, ERR, and DQ arrays are
consumed directly, and no frozen background subtraction is applied to the
science data.

The unit of the NOVA-S likelihood is one retained physical detector pixel in
one integration. Valid samples have finite values, positive uncertainties,
and no \texttt{DO\_NOT\_USE} flag; the retained detector matrix contains
97{,}557 pixels per integration. The detector support is organised into
native (order, detector-column) groups. Six edge groups without reliable
mapping are excluded, leaving 2{,}720 native groups. The two orders retain
disjoint support in the current release, and predicted order contributions
are nevertheless summed in detector space before each residual is formed, so
the formulation extends unchanged to overlapping support.

\subsection{The \nova{}-S detector-level forward model}

% Question: What does NOVA-S predict for each detector sample?
% Answer: The sum over orders of a response-propagated product of the frozen
% source-curvature factor, the transit, and native-column continua, plus an
% additive background.
% Evidence: The exact NOVA-S method note, Section 3.

For integration $t$ and physical detector pixel $p$, NOVA-S predicts
\begin{equation}
\begin{split}
 \widehat y_{tp}
 = {}& \sum_o \left[\bm{R}_o\left\{
   g(t)\,\bm{T}_{o,t}(D,\bm{u}_o)\odot
   \bm{\Psi}_o\,\bm{c}_o(t)\right\}\right]_p \\
 & + \left[\bm{H}\bm{\alpha}\right]_{tp} ,
\end{split}
 \label{eq:nova-forward}
\end{equation}
where $o$ indexes the spectral orders. The frozen response $\bm{R}_o$
contains the PASTASOSS wavelength and trace geometry
\citep{BainesEtAl2023Trace,BainesEtAl2023Wavelength}, the WebbKernel
line-spread mapping distributed with ATOCA \citep{DarveauBernierEtAl2022},
and the empirical q-star cross-dispersion redistribution described below.
$\bm{\Psi}_o$ maps native-column amplitudes onto the fine wavelength
representation, $\bm{T}_{o,t}$ is the exact exposure-integrated transit for
depth spectrum $D$ and limb darkening $\bm{u}_o$, $g(t)$ is a frozen common
source-curvature factor, $\bm{c}_o(t)$ collects the native per-column
stellar continua, and $\bm{H}\bm{\alpha}$ is the additive E13b background.
The symbol $\odot$ denotes element-wise multiplication on the wavelength
grid. The source and transit terms of both orders are added before the
residual is formed. The background is additive and is never multiplied by
the transit or by $g(t)$: a change in transit depth cannot be compensated by
attenuating the background in the same way as the star.

\subsection{Shared transmission spectrum and transit model}

% Question: How is the transmission spectrum parameterised, and what orbital
% and limb-darkening model generates the transit?
% Answer: 160 shared positive adaptive coefficients split into an achromatic
% depth and a weighted-zero-mean morphology; fixed Keplerian geometry with
% finite-exposure quadrature; eight fitted limb-darkening coordinates around
% the STAGGER theory reference.
% Evidence: The exact NOVA-S method note, Sections 4 and 5.

The transmission spectrum is represented by $K=160$ shared adaptive physical
basis coefficients,
\begin{equation}
  D_j = D_{\mathrm{global}}\,
  \frac{\exp\!\left[(\bm{H}_v\bm{v})_j\right]}
       {\sum_k w_k \exp\!\left[(\bm{H}_v\bm{v})_k\right]} ,
  \label{eq:depth-basis}
\end{equation}
where $\bm{v}$ is a signed bounded optimiser coordinate, $w_k$ are frozen
basis-width weights, and $\bm{H}_v$ removes the constant gauge. This
separates the achromatic depth $D_{\mathrm{global}}$, which carries a small
positive floor, from a weighted-zero-mean chromatic morphology, and it
enforces positive depths. Orders 1 and 2 evaluate the same physical
$D(\lambda)$ on their own wavelength grids. Two optional order-discrepancy
coordinates exist in the code but are numerically fixed to zero within
$\pm 10^{-12}$; they have no physical freedom in NOVA-S. No atmospheric
template, wavelength-local smoother, or truth-informed spectral prior is
used.

The orbital geometry is fixed: period, scaled semi-major axis, impact
parameter, and mid-transit time are held at the values listed in
Appendix~\ref{app:nova}, with zero eccentricity. Transit light curves are
evaluated with an exact Keplerian engine implemented in jaxoplanet
\citep{Jaxoplanet2024}, and each detector integration is modelled as a
finite-exposure average over 21 quadrature points rather than at a single
mid-exposure time.

Quadratic limb darkening \citep{MandelAgol2002} is parameterised through
bounded coordinates \citep{Kipping2013},
\begin{equation}
 u_1=2\sqrt{q_1}\,q_2, \qquad
 u_2=\sqrt{q_1}\,(1-2q_2).
 \label{eq:kipping-coordinates}
\end{equation}
The wavelength-dependent theory reference is the STAGGER-grid prediction
\citep{MagicEtAl2015} and is unchanged in NOVA-S; the fitted freedom around
it is deliberately small. For each order, six knots describe each of $q_1$
and $q_2$, but only a constant offset and a normalised log-wavelength slope
are fitted for each variable and order, giving eight nonlinear
limb-darkening coordinates in total, projected into the six-knot deviation
basis. Profiled priors constrain the deviation from theory and the
wavelength curvature, and fine nodes outside the valid theory support are
excluded rather than extrapolated.

\subsection{Response, native continua, uncertainties, and background}

% Question: Which frozen components close the model between transit and
% detector, and how were they calibrated without truth access?
% Answer: An OOT-estimated q-star cross-dispersion response with documented
% limitations, 5,440 native continuum coefficients, per-order and per-row
% uncertainty scales from OOT cross-validation, and the 16-coefficient E13b
% additive background anchored on off-trace pixels.
% Evidence: The exact NOVA-S method note, Sections 6 to 9.

Within each native group, the empirical q-star response redistributes the
group flux across detector rows; it is a normalised spatial response, not an
additional spectrum. It was rebuilt from the public V45 cube itself, without
access to injected truth or generator-private products: pixel time series
were fitted with a level, a linear trend, and a depth-free event morphology
under robust loss, alternating phase-stratified folds provided independent
estimates, and the full response was admitted only after cross-fold
flux-distribution and centroid stability gates. The frozen response is
retained with a documented limitation: its smoother operates on an integer
trace-relative coordinate, which produces centroid resets and a small
red-column amplitude mismatch. The construction, admission gates, and this
defect are specified in Appendix~\ref{app:nova}; its causal assessment
belongs to the Results.

Each of the 2{,}720 native groups carries a continuum level and a linear
time slope, giving 5{,}440 linear nuisance coefficients that are eliminated
exactly at every nonlinear evaluation. The group continuum controls total
group brightness, while q-star controls only the within-group spatial
distribution; this separation is why a response-amplitude error relative to
the delivered source can act multiplicatively on the recovered depth.

Detector uncertainties start from the FITS ERR arrays, scaled by one frozen
factor per order, and are then multiplied by a per-detector-row guard
derived from bidirectional OOT cross-validation, floored at one. No
wavelength information, injected truth, or recovered spectrum enters this
rule; the exact scales and formula are given in Appendix~\ref{app:nova}.

The E13b additive background contains eight spatial modes crossed with two
temporal modes, giving 16 linear coefficients. The spatial family and rank
were selected by blocked cross-validation on authenticated off-trace pixels,
and the second temporal mode was selected from the leading OOT singular
vectors subject to a fixed cap on its absolute correlation with the transit
reference vector. The OOT expectation and normal matrix of these
coefficients form a quadratic anchor, and background and continua are fitted
jointly in one linear solve. There is no frozen background subtraction in
recovery. The likelihood contains no generic Gaussian-process term for
time-correlated residuals; the robust weights described below reduce the
influence of outliers, and residual temporal correlation is retained as a
diagnostic.

\subsection{Out-of-transit common source curvature}

% Question: How is slow temporal variation of the source handled?
% Answer: A frozen wavelength-independent curvature factor calibrated on
% held-out OOT folds; the common candidate was selected because the per-order
% candidate failed the requirement that both orders improve.
% Evidence: The exact NOVA-S method note, Section 10.

A small wavelength-independent temporal source-curvature factor $g(t)$ was
calibrated from the 102 complement integrations, split into two 51-integration
folds. After subtracting only the E13b OOT expectation for this calibration
step, stable trace flux was integrated per order, weighted quadratic and
affine time models were compared, and their positive ratio defined a
candidate curvature. Null, common, and per-order candidates were ranked by
held-out standardised chi-squared. The per-order candidate had the lowest
aggregate score but failed the preregistered requirement that both orders
improve individually, so the common candidate was selected. The frozen
factor has unit OOT mean, multiplies only the stellar and transit source
bundle, never the background, and adds no fitted parameters. Its selection
scores and range are listed in Appendix~\ref{app:nova}.

\subsection{Profiled robust optimisation and convergence}

% Question: How is the model fitted, and what counts as convergence?
% Answer: Exact variable projection of all 5,456 linear coefficients inside
% bounded trust-region least squares, wrapped in exact outer Huber IRLS, with
% preregistered convergence gates and two independent deterministic starts.
% Evidence: The exact NOVA-S method note, Section 11.

Let $\bm{\theta}$ contain the nonlinear spectrum and limb-darkening
parameters and $\bm{\eta}$ the linear continuum and background coefficients.
At fixed $\bm{\theta}$ and fixed robust weights $\bm{W}$, NOVA-S solves the
single global linear problem
\begin{equation}
\begin{split}
 \widehat{\bm{\eta}}(\bm{\theta})=
 \arg\min_{\bm{\eta}}\;
 &\frac{1}{2}\left\|\bm{W}^{1/2}
 [\bm{y}-\bm{A}(\bm{\theta})\bm{\eta}]\right\|_2^2 \\
 &+\frac{1}{2}
 \left\|\bm{L}(\bm{\alpha}-\bm{\alpha}_{\mathrm{off}})\right\|_2^2 ,
\end{split}
 \label{eq:varpro}
\end{equation}
where the regularisation row anchors the background on its OOT expectation
with $\bm{L}^{\mathsf T}\bm{L}$ equal to the OOT normal matrix. This
variable-projection step \citep{GolubPereyra1973} eliminates all 5{,}440
continuum and 16 background coefficients exactly: the equilibrated sparse
continuum block is factorised and the background coupling is reduced to a
16-dimensional Schur complement. The profiled Jacobian differentiates
through the exact inner optimum and is accumulated with streamed sufficient
statistics, avoiding a visit-sized dense Jacobian.

The outer nonlinear problem is solved with bounded trust-region reflective
least squares \citep{BranchColemanLi1999}. Robustness is implemented as
exact outer Huber iteratively reweighted least squares \citep{Huber1964}
with threshold 1.345, so the inner solver retains a linear least-squares
loss and the combination remains exact M-estimation. Convergence requires
robust-weight stability, small relative objective change, small
shared-spectrum changes, scale-aware Karush-Kuhn-Tucker and first-order
criteria, agreement of the deterministic starts, and a final confirmation
fit at the recomputed terminal weights; the numeric gates and cycle caps are
listed in Appendix~\ref{app:nova}. Two deterministic starts are fitted
independently with no state transfer: one at the frozen white-light depth
with zero morphology and theory limb darkening, and one with a fixed
zero-mean 250 ppm morphology perturbation and small fixed limb-darkening
perturbations. Both starts of the frozen release converged at robust cycle
14. The selected solution is determined by the prospective convergence
policy, never by truth error.

\subsection{Benchmark separation, comparison analyses, and fairness status}

% Question: How are injector, recovery, and comparison methods kept
% independent, and what fairness work remains open?
% Answer: Generator, delivery, recovery, and scoring are isolated; one opaque
% cube is consumed unchanged by every pipeline; an authentic Ahsoka analysis
% is the first comparison; a paired group-stage and rateints-stage injection
% experiment is preregistered but not complete.
% Evidence: The injector evolution and fairness record.

Injector and recovery products are separated by construction. Recovery code
reads only the delivered cube and preregistered auxiliary products; it
cannot read the atmospheric spectrum, generator-only response products, or
scoring code. Hashes identify the delivered arrays, configuration, masks,
and response products, and the injected spectrum is loaded only after a
prediction and its provenance record have been closed
(Appendix~\ref{app:benchmarking}). The injector consumes none of NOVA's
q-star, masks, recovered spectra, or truth residuals; its source deficit is
built independently per order from the same public instrument references.

The first external comparison uses an authentic Ahsoka reduction of the same
delivered cube, following the analysis topology of the published WASP-17~b
spectrum, including detector preparation, spectral extraction, spectroscopic
light curves, and channel-level transit inference \citep{LouieEtAl2025}.
Fairness is defined at delivery rather than by forcing a common internal
representation: pixel values, uncertainties, flags, time grid, and declared
wavelength support are identical, unsupported bins are excluded
symmetrically, and each method retains its intended nuisance model and
inference strategy.

One fairness limitation is open. The principal comparison pipelines perform
their most distinctive $1/f$ correction on detector groups before ramp
fitting, and a \texttt{rateints} cube contains no groups. The programme
therefore includes a preregistered paired experiment that injects the same
physical deficit at group stage and at \texttt{rateints} stage, runs NOVA
and the comparison pipelines on both, and compares delivered deficits,
white-light scatter, and common-bin recovery. This is a sensitivity
experiment for the frozen V45 control; it has not been completed and is not
part of the frozen result. Later benchmark rounds will include Eureka!,
exoTEDRF, transitspectroscopy, and the official JWST pipeline with ATOCA
extraction
\citep{BellEtAl2022,Radica2024,DarveauBernierEtAl2022,LouieEtAl2025},
reported both end-to-end and with a common downstream light-curve fitter;
none of these pipelines has yet been run with its authentic group-level
preprocessing on this benchmark.

\subsection{Reporting support and evaluation metrics}

% Question: On what wavelength support are methods compared, and with which
% metrics?
% Answer: A common R=100 grid restricted prospectively to the direct
% PASTASOSS training domain, scored by absolute RMSE, mean bias, and
% morphology RMSE, alongside uncertainty and validity diagnostics.
% Evidence: The exact NOVA-S method note and appendix draft.

Recovered and injected spectra are compared on a common $R=100$ reporting
grid. One instrument-defined rule fixes the paper-wide support: any bin
whose detector support extends beyond the direct PASTASOSS training domain
is excluded in full, because the wavelength solution beyond the last
training anchor is a polynomial extrapolation. For this visit the rule
retains 141 bins, with the last whole bin ending near $2.740\,\mu\mathrm{m}$.
The rule is fixed across methods, changes no fit or prediction, and is
applied only after recovery; its exact detector-coordinate definition is
given in Appendix~\ref{app:benchmarking}. Internal regression grids are
retained for diagnostics.

Let $\widehat D_i$ and $D_i^{\mathrm{true}}$ be the recovered and injected
depths on the common grid and $e_i=\widehat D_i-D_i^{\mathrm{true}}$. The
absolute recovery error is
\begin{equation}
 \mathrm{RMSE}_{\mathrm{abs}}=
 \left(\frac{1}{N}\sum_i e_i^2\right)^{1/2},
 \label{eq:absolute-rmse}
\end{equation}
and the morphology error removes the mean residual,
\begin{equation}
 \mathrm{RMSE}_{\mathrm{morph}}=
 \left[\frac{1}{N}\sum_i(e_i-\bar e)^2\right]^{1/2},
 \label{eq:morphology-rmse}
\end{equation}
so that a constant offset is distinguished from incorrect spectral shape.
Mean and median bias, feature-amplitude error, order-specific residuals, and
subregion morphology are reported alongside. Uncertainty calibration is
evaluated through standardised residuals, empirical coverage of injected
depths, and repeated noise realisations, with spectral covariance from the
profiled normal system where available. Detector residuals, robust weights,
rank, conditioning, boundary activity, and convergence diagnostics accompany
every score: a low RMSE is not sufficient if a fit fails its preregistered
numerical or physical validity gates.
